% ==========================================================
% PREÁMBULO: paquetes, colores y comandos reutilizables
% Edita este archivo para cambiar tipografía, márgenes,
% colores de los badges, o el formato de los títulos.
% ==========================================================

\usepackage[top=3cm,left=4cm,bottom=3cm,right=2cm]{geometry}
\usepackage[utf8]{inputenc}
\usepackage[T1]{fontenc}
\usepackage{mathptmx} % tipografía tipo Times, similar al PDF de referencia
\usepackage{setspace}
\usepackage{titlesec}
\usepackage{tabularx}
\usepackage{longtable}
\usepackage{booktabs}
\usepackage{array}
\usepackage[table]{xcolor}
\usepackage[most]{tcolorbox}
\usepackage{enumitem}
\usepackage{hyperref}
\hypersetup{colorlinks=false, hidelinks}
\usepackage{xurl}

\onehalfspacing

% ---------- Colores (tono ámbar del badge "EN CONSTRUCCIÓN") ----------
\definecolor{badgebg}{HTML}{FDF6E3}
\definecolor{badgeborder}{HTML}{C9A227}
\definecolor{badgetext}{HTML}{8A6D1D}
\definecolor{greensignal}{HTML}{2E7D32}

% ---------- Formato de títulos de sección (centrado, versalitas, como el PDF) ----------
\titleformat{\section}{\centering\bfseries\large}{\thesection.}{0.5em}{\MakeUppercase}
\titleformat{\subsection}{\bfseries\normalsize}{\thesubsection}{0.5em}{}
\titlespacing*{\section}{0pt}{2.5em}{1.5em}

% ---------- Comando para el "badge" EN CONSTRUCCIÓN / estado ----------
% Uso: \estado{En construcción}
\newcommand{\estado}[1]{%
  \begin{center}
  \tcbox[colback=badgebg,colframe=badgeborder,boxrule=0.5pt,
    arc=6pt,left=6pt,right=6pt,top=3pt,bottom=3pt,
    fontupper=\small\bfseries\color{badgetext}]{\MakeUppercase{#1}}
  \end{center}
}

% ---------- Comando para la nota itálica que acompaña cada estado ----------
% Uso: \notafase{Se redacta en la Fase 3 del curso.}
\newcommand{\notafase}[1]{\begin{center}\itshape\small #1\end{center}}

% ---------- Columna de tabla con ancho fijo y texto justificado a la izquierda ----------
\newcolumntype{L}[1]{>{\raggedright\arraybackslash}p{#1}}
